% pxPM -- P3CTeX project management
\documentclass[a4paper]{article}
\usepackage[english]{babel}
\usepackage{pxPM}
\begin{document}
\section*{pxPM: project management}
\section{Daily Gantt charts}
The \texttt{pxPM} planning tool provides a three-command interface to
\texttt{pgfgantt}; it does not change the document orientation.
Select it when loading P3CTeX: \verb|\documentclass[project-management,PR,CAT]{P3CTeX}|.
Outside the class, use \verb|\usepackage{pxPM}|.

\begin{verbatim}
\pxGantt{start=01/10/2026,end=31/01/2027,days-per-block=28}
\pxGanttItem{Investigation}{01/10/2026}{20/10/2026}
\pxGanttItem{Implementation}{15/10/2026}{15/12/2026}
\pxGanttEnd
\end{verbatim}

Options: \texttt{start}, \texttt{end} (required dd/mm/yyyy),
\texttt{days-per-block} (default 28), \texttt{days-per-page} (alias),
\texttt{day-width} (optional TeX length), \texttt{label-width} (manual override),
\texttt{label-max-width} (default \verb|.24\linewidth|),
\texttt{rotate=true} (optional chart-only rotation).

By default the label column is measured from actual task text and capped
relative to the available line width. Remaining width is divided among
calendar days. The date range is balanced into consecutive blocks, and
bars crossing block boundaries are clipped. P3CTeX colours are reused;
standalone users receive fallback colours.

\section*{Example}
\pxGantt{start=01/10/2026,end=25/11/2026,days-per-block=28}
\pxGanttItem{Investigació}{01/10/2026}{20/10/2026}
\pxGanttItem{Disseny}{15/10/2026}{10/11/2026}
\pxGanttItem{Implementació}{01/11/2026}{25/11/2026}
\pxGanttEnd
\section{Project planning tools}
Load \texttt{project-management} in the P3CTeX class, or load \texttt{pxPM}
with a standard document class. The same tool provides SWOT/DAFO tables,
user stories, scenarios, use-case descriptions and Gantt schedules.
Labels follow P3CTeX's Catalan, Spanish and English language selection.
\begin{verbatim}
\pxSWOT{Strengths}{Weaknesses}{Opportunities}{Threats}
\begin{pxUserStory}{Reserve equipment}
  \pxAsA{a customer}
  \pxIWant{to reserve equipment}
  \pxSoThat{I can plan a shoot}
  \pxAcceptanceCriteria{The reservation is confirmed.}
\end{pxUserStory}
\begin{pxScenario}{Successful reservation}
  \pxGiven{available equipment}
  \pxWhen{a customer confirms}
  \pxThen{a reservation is created}
\end{pxScenario}
\begin{pxUseCase}{Reserve equipment}
  \pxActors{Customer}
  \pxPreconditions{Equipment is available.}
  \pxMainFlow{Select equipment, then confirm.}
  \pxAlternatives{Choose another date if unavailable.}
  \pxPostconditions{A reservation exists.}
\end{pxUseCase}
\end{verbatim}
Fields accept paragraphs and shared preamble commands. Card environments accept
local tcolorbox options; \verb|\pxPMsetup{labels={...},box-options={...}}|
customises defaults. Gantt commands below belong to the same package.
\end{document}
